% Generated draft: requires mathematical and source review.
\documentclass[11pt]{article}
\usepackage[margin=25mm]{geometry}
\usepackage{fontspec}
\setmainfont{Latin Modern Roman}
\usepackage{amsmath,amssymb,mathtools,unicode-math}
\setmathfont{Latin Modern Math}
\usepackage{enumitem,needspace,xurl,hyperref}
\setlength{\parindent}{0pt}
\setlength{\parskip}{5pt}
\setlength{\emergencystretch}{3em}
\newcommand{\N}{\mathbb{N}}
\newcommand{\Z}{\mathbb{Z}}
\newcommand{\Q}{\mathbb{Q}}
\newcommand{\R}{\mathbb{R}}
\newcommand{\C}{\mathbb{C}}
\newcommand{\F}{\mathbb{F}}
\newcommand{\A}{\mathbb{A}}
\newcommand{\PP}{\mathbb{P}}
\newcommand{\E}{\mathbb{E}}
\newcommand{\eps}{\varepsilon}
\newcommand{\dd}{\,\mathrm{d}}
\begin{document}
\section*{UnsolvedMath ID: 700917}
\textbf{Source problem code:} catalogue entry
\subsection{Definitions and mathematical statement}
An involution in a group is an element of order exactly \(2\). Write \(x^g=g^{-1}xg\), \(i^G=\{i^g:g\in G\}\), and \(o(x)\) for the order of an element \(x\). For a subgroup \(H\leq G\), put
\[
N_G(H)=\{g\in G:g^{-1}Hg=H\}.
\]
A subgroup \(H\) is \(i\)-invariant if \(H^i=H\). A \(2\)-subgroup is understood here to mean a subgroup every element of which has finite order a power of \(2\); the source does not explicitly restrict these subgroups to be finite. A simple group is a nontrivial group whose only normal subgroups are itself and the identity subgroup.

\textbf{Unresolved terminology.} The supplied record does not define the \emph{periodic part} of a group. Write \(P(H)\) for that source term, without identifying it with either the set of finite-order elements or a particular periodic subgroup. These interpretations require confirmation, since the finite-order elements of an arbitrary group need not form a subgroup. Likewise, let \(\mathcal Q\) denote the source's intended class of generalized quaternion groups. The usual finite generalized quaternion groups have presentations
\[
Q_{2^n}=\langle a,b\mid a^{2^{n-1}}=1,\ b^2=a^{2^{n-2}},\ b^{-1}ab=a^{-1}\rangle,
\qquad n\geq3.
\]
The record does not specify whether an infinite generalized quaternion group is also permitted. Finally, let \(\mathcal F\) denote the source's class of Frobenius groups with abelian kernel and finite complement of even order. Its referenced definition in Problem 6.55 is not supplied; in particular, the intended convention for possibly infinite Frobenius groups remains unresolved.

\textbf{Source formulation, conditional on those terminology conventions.} A group \(G\), together with a distinguished involution \(i\in G\), satisfies the stated \(T_0\) conditions when, on putting
\[
C=C_G(i)=\{g\in G:gi=ig\},
\]
the following hold:
\begin{enumerate}
\item For every \(u,v\in i^G\), one has \(o(uv)<\infty\).
\item Every \(2\)-subgroup \(S\leq G\) is cyclic or belongs to \(\mathcal Q\).
\item The subgroup \(C\) is infinite, \(C\ne G\), and its periodic part is finite: \(|P(C)|<\infty\).
\item For every subgroup \(H\leq G\) satisfying \(1<|H|<\infty\) and \(H^i=H\),
\[
N_G(H)\leq C
\quad\text{or}\quad
P\bigl(N_G(H)\bigr)\in\mathcal F.
\]
In the second alternative, the periodic part is asserted to be a group with the specified Frobenius structure.
\item For every \(c\in G\setminus C\) such that \(o(ci)=2\), there exists \(s\in C\) such that
\[
\bigl|\langle c,c^s\rangle\bigr|=\infty.
\]
\end{enumerate}
The question is whether there exist a simple group \(G\) and an involution \(i\in G\) satisfying all five conditions. The distinguished involution must be the same throughout. This is a statement-only transcription with explicitly unresolved terminology, rather than a fully disambiguated definition of \(T_0\).
\subsection{Short English statement}
Does there exist a simple group containing an involution for which all five stated \(T_0\) conditions hold? These conditions constrain products of conjugate involutions, \(2\)-subgroups, the involution's centralizer, certain finite-subgroup normalizers, and subgroups generated by an element and one of its centralizer conjugates. The meanings of periodic part, generalized quaternion in the infinite setting, and the referenced Frobenius-group convention require confirmation.
\subsection{Sources}
\begin{itemize}
\item Local UnsolvedMath registry, permanent ID \(700917\), source code \texttt{catalogue entry}, title \emph{Kourovka Notebook Problem 12.101}. The supplied record labels the problem open; that label is supplied registry metadata, not an independently verified literature assessment.
\item V. P. Shunkov, proposer, 1992. Attributed source: \emph{The Kourovka Notebook (Unsolved Problems in Group Theory)}, 12th issue (1992), Problem 12.101; supplied edition reference: 21st issue edition (2026), \texttt{arXiv:1401.0300v48}. The collection description names E. I. Khukhro and V. D. Mazurov as editors.
\item The statement refers to Kourovka Notebook Problem 6.55 for Frobenius groups. That referenced text is not included in the supplied record.
\end{itemize}
\end{document}
